% !TeX spellcheck = de_DE

Die Definition und Strukturierung von Anforderungen bildet die Grundlage für jedes erfolgreiche \ac{it}-Projekt. Sie dienen als gemeinsame Verständigungsbasis zwischen Auftraggeber und Auftragnehmer und schaffen Klarheit über Ziele, Rahmenbedingungen und erwartete Funktionalitäten. Studien zeigen, dass Fehler im Umgang mit Anforderungen erhebliche Kostenverursacher in Projekten darstellen und damit maßgeblich den Projekterfolg gefährden können~\cite[94]{ruppRequirementsEngineeringUnd2009}. Um solchen Risiken vorzubeugen, empfehlen internationale Standards wie \ac{ieee} 830~\cite{IEEERecommendedPractice1998} eine vollständige, eindeutige und nachvollziehbare Dokumentation, da die Qualität eines Konzepts wesentlich von der Qualität der zugrunde liegenden Anforderungen abhängt~\cite[1f]{chikhNewTraceableSoftware2012}. Für die Entwicklung eines \ac{sd-wan}-Konzepts ist dies von besonderer Bedeutung, da hierbei sowohl technische als auch organisatorische und sicherheitsrelevante Anforderungen präzise erfasst und strukturiert werden müssen. Vor diesem Hintergrund werden im Folgenden die spezifischen Anforderungen an das zu entwickelnde Konzept dargestellt. \\
Diese teilen sich in funktionale und nicht-funktionale Anforderungen auf. Diese Struktur entspricht der Idee von \ac{ieee} 830, Anforderungen systematisch nach Funktionalität zu klassifizieren, um die Übersichtlichkeit zu gewährleisten~\cite{chikhNewTraceableSoftware2012}. \\
\\
Die folgenden Anforderungen wurden in enger Zusammenarbeit mit den jeweiligen Fachbereichen erarbeitet. Grundlage hierfür waren Gespräche und Interviews mit den fachlich verantwortlichen Personen, um die tatsächlichen Bedarfe und Herausforderungen aus der Praxis zu erfassen. Die Priorisierung der Anforderungen erfolgte ebenfalls auf Basis des Expertenwissens und der Einschätzungen der jeweiligen Fachbereichsvertreter. Dadurch wird sichergestellt, dass die Bewertung der Anforderungen die tatsächlichen betrieblichen Notwendigkeiten und das vorhandene Know-how widerspiegelt.

\subsubsection*{Funktionale Anforderungen}
Funktionale Anforderungen beschreiben die zentralen Eigenschaften und Fähigkeiten, die eine \ac{sd-wan}-Lösung grundsätzlich bereitstellen muss. Diese wurden in \ref{sec:sd-wan-grundlagen} bereits beschrieben. Es wird davon ausgegangen, dass das \ac{sd-wan}-System diese Funktionen bereitstellt. Diese funktionalen Anforderungen beziehen sich damit weniger auf das konkrete Konzept dieser Arbeit, sondern auf die generellen Funktionsmerkmale der Technologie~\cite[207]{broyProduktUndQualitaetsanforderungen2021}. Solche Anforderungen umfassen beispielsweise die Möglichkeit zur dynamischen Pfadwahl über verschiedene Transportnetze, die zentrale Verwaltung sowie die anwendungsorientierte Steuerung von Datenflüssen. (Kapitel \ref{sec:sd-wan-grundlagen}) Diese Funktionen bilden den technischen Rahmen, auf dessen Basis im weiteren Verlauf die spezifischen, unternehmensbezogenen Anforderungen an das zu entwickelnde \ac{sd-wan}-Konzept abgeleitet werden.

\subsubsection*{Nicht-Funktionale Anforderungen}

\enquote{Nicht-Funktionale Anforderungen beschreiben Anforderungen an ein System und Randbedingungen, welche die Art der Erbringung der in den funktionalen Anforderungen festgelegten Funktionen beschreiben.}~\cite[207]{broyProduktUndQualitaetsanforderungen2021}\\
Die nicht-funktionalen Anforderungen sind in Form einer Tabelle zusammengefasst und nach \ac{ieee} spezifiziert.

\begin{table}[H]
	\centering
	\begin{tabularx}{\textwidth}{lXll}
		\hline
		\textbf{Nr.} & \textbf{Beschreibung} & \textbf{Priorität} & \textbf{Kategorie} \\
		\hline
		ANF-01 & Das \ac{sd-wan} muss eine stabile und zuverlässige Anbindung aller Standorte gewährleisten. & \cellcolor{red!30}Hoch & Konnektivität \\
		ANF-02 & Das \ac{sd-wan} muss Local Internet Breakout unterstützen, um Internet-Traffic direkt vor Ort ausleiten zu können. & \cellcolor{red!30}Hoch & Performance \\
		ANF-03 & Das \ac{sd-wan} soll die Bündelung mehrerer \ac{wan}-Leitungen ermöglichen, um Ausfallsicherheit und Bandbreite zu erhöhen. & \cellcolor{red!30}Hoch & Redundanz \\
		ANF-04 & Das \ac{sd-wan} muss mit Leitungen unterschiedlicher Qualität umgehen können und dynamisch anpassen. & \cellcolor{red!30}Hoch & Flexibilität \\
		ANF-05 & Das Konzept muss Video-Streaming mit hoher Bandbreite und niedriger Latenz unterstützen. & \cellcolor{yellow!30}Mittel & Performance \\
		ANF-06 & Das \ac{sd-wan}-Konzept muss eine sichere Verbindung zum Hauptstandort gewährleisten. & \cellcolor{red!30}Hoch & Sicherheit \\
		ANF-07 & Das Konzept muss zuverlässig und ohne größeren administrativen Aufwand funktionieren. & \cellcolor{yellow!30}Mittel & Betrieb \\
		\hline
	\end{tabularx}
	\caption{Anforderungen an die \ac{sd-wan}-Lösung}
\end{table}

Das zentrale Ziel des \ac{sd-wan}-Konzepts besteht darin, die stabile und zuverlässige Anbindung der Auslandsstudios (ANF-01) sicherzustellen, unabhängig von der Art der verwendeten Leitungen~\cite{heckmannInterviewMitSystemtechnik2025}. Wie in Abschnitt~\ref{sec:ausgangslage} beschrieben, geht es hierbei insbesondere darum, die bisherige Abhängigkeit von festen Layer-2-Verbindungen aufzulösen.\\
Der Local Internet Breakout ist von zentraler Bedeutung, da er es den Standorten ermöglicht, Internet-Traffic direkt vor Ort auszuliefern. Mit der bisherigen Layer-2-Konfiguration konnte diese Funktion nicht realisiert werden, sodass sämtlicher Internet-Traffic über die zentralen Standorte in Deutschland geleitet wird. Dies führt dazu, dass Mitarbeiter der Auslandsstudios teilweise keinen direkten Zugriff auf lokale Internetdienste, etwa Nachrichtenseiten oder branchenspezifische Portale, haben. Die eingeschränkte Verfügbarkeit lokaler Internetinhalte beeinträchtigt die Arbeit erheblich, insbesondere bei zeitkritischen Recherchen und der Nutzung standortspezifischer Anwendungen. Daher ist die Implementierung eines Local Internet Breakout ein wesentlicher Bestandteil des Konzepts~\cite{heckmannInterviewMitSystemtechnik2025}.\\
Ein weiterer zentraler Aspekt ist der Umgang mit Leitungen unterschiedlicher Qualität, da die Bandbreite, Latenz und Zuverlässigkeit stark von den jeweiligen Anbietern und der geografischen Lage der Standorte abhängen. So haben manche Standorte nur eine gute Layer-3 Verbindung und mehrere zusätzliche Layer-3 Verbindungen mit geringerer Bandbreite und Verbindungsqualität. Das Konzept muss daher sicherstellen, dass die Lösung trotz unterschiedlicher Leitungsqualitäten stabile Verbindungen, ausreichende Bandbreite für geschäftskritische Anwendungen und eine konsistente Servicequalität gewährleistet.\\
Für den \ac{swr} ist Video-Streaming aus den verschiedenen Standorten ein zentraler Bestandteil des täglichen Betriebs. Die Übertragung dieser Streams muss zuverlässig, mit hoher Bandbreite und niedriger Latenz erfolgen, da die Produktion und Verteilung von Live-Inhalten unmittelbar von der stabilen Netzwerkverbindung abhängt. Aus diesem Grund ist das \ac{sd-wan} so zu konzipieren, dass es diese Anforderungen für alle Standorte gewährleistet und Ausfälle oder Verzögerungen im Stream vermeidet.\\
Für das Konzept sind vor allem die Anforderungen relevant, welche eine hohe Priorität haben. Diese müssen für die erfolgreiche Erstellung des Konzepts unbedingt berücksichtigt werden. Eine Nicht-Erfüllung einer dieser Anforderungen macht das Konzept nicht einsetzbar.

\subsubsection*{Risiken bei Nichterfüllung}
Nachdem im vorhergehenden Kapitel die Anforderungen an das zu entwickelnde \ac{sd-wan}-Konzept detailliert dargestellt wurden, ist es entscheidend, mögliche Risiken zu identifizieren, die durch eine Nichterfüllung dieser Anforderungen entstehen können.

\begin{itemize}
	\item \textbf{ANF-01:} Ausfälle oder Unterbrechungen im Netzwerk, verzögerte Produktion, Beeinträchtigung des laufenden Betriebs, erhöhte Support- und Betriebskosten
	\item \textbf{ANF-02:} Internet-Traffic wird weiterhin über zentrale Standorte geleitet $\rightarrow$ erhöhte Latenz, erhöhter Traffic auf \ac{vpn}, eingeschränkter Zugriff auf lokale Inhalte, reduzierte Produktivität
	\item \textbf{ANF-03:} Reduzierte Bandbreite und Ausfallsicherheit, Single-Point of Failure, höhere Wahrscheinlichkeit von Netzwerküberlastungen, eingeschränkte Skalierbarkeit
	\item \textbf{ANF-04:} Instabile Verbindungen, Paketverlust, ungleichmäßige Performance, Unterbrechungen kritischer Anwendungen, inkonsistente Servicequalität zwischen Standorten
	\item \textbf{ANF-05:} Abbrüche oder Verzögerungen bei Live-Streams, schlechtere Qualität der Medienproduktion, Unzufriedenheit der Zuschauer und Nutzer, mögliche Produktionsausfälle
	\item \textbf{ANF-06:} Sicherheitslücken, erhöhte Angriffsfläche, potenzieller Datenverlust, erhöhte Risiken für Unternehmensinformationen
	\item \textbf{ANF-07:} Hoher administrativer Aufwand, fehleranfällige Konfigurationen, langsame Problemlösung, Betriebskosten steigen, geringere Akzeptanz durch Mitarbeiter
\end{itemize}